\ifdefined\gkver\else\def\gkver{student}\fi
\documentclass[\gkver]{gaokaozhenti}
\begin{document}

\chapter{2023年全国新课标}
%% paperType: 理综物理部分

\begin{enumerate}
\item
%% number: 1
%% paperName: 2023年全国高考新课标第 1 题 6 分
%% typeId: 1
%% type: 单选题
%% chapter: 机械波
%% point: 波长频率波速
%% method: 无
%% score: 6
%% degree: 950
%% duplicateId: 0
%% body:
船上的人和水下的潜水员都能听见轮船的鸣笛声。声波在空气中和在水中传播时的\xzanswer{A}

\fourchoices[answer=A]
{波速和波长均不同}
{频率和波速均不同}
{波长和周期均不同}
{周期和频率均不同}

%% answer: A
%% memo:
\memoanswer{
声波的周期和频率由振源决定，故声波在空气中和在水中传播的周期和频率均相同，但声波在空气和水中传播的波速不同，根据波速与波长关系 $v = \lambda f$ 可知，波长也不同。故 A 正确，BCD 错误。

故选 A。
}

\item
%% number: 2
%% paperName: 2023年全国高考新课标第 2 题 6 分
%% typeId: 1
%% type: 单选题
%% chapter: 功和能
%% point: 动能定理
%% method: 无
%% score: 6
%% degree: 950
%% duplicateId: 0
%% body:
无风时，雨滴受空气阻力的作用在地面附近会以恒定的速率竖直下落。一质量为 $m$ 的雨滴在地面附近以速率 $v$ 下落高度 $h$ 的过程中，克服空气阻力做的功为（重力加速度大小为 $g$）\xzanswer{B}

\fourchoices[answer=B]
{0}
{$mgh$}
{$\frac{1}{2}mv^{2} - mgh$}
{$\frac{1}{2}mv^{2} + mgh$}

%% answer: B
%% memo:
\memoanswer{
在地面附近雨滴做匀速运动，根据动能定理得

$mgh - W_{f} = 0$

故雨滴克服空气阻力做功为 $mgh$。

故选 B。
}

\item
%% number: 3
%% paperName: 2023年全国高考新课标第 3 题 6 分
%% typeId: 1
%% type: 单选题
%% chapter: 光学
%% point: 光子说
%% method: 无
%% score: 6
%% degree: 850
%% duplicateId: 0
%% body:
铯原子基态的两个超精细能级之间跃迁发射的光子具有稳定的频率，铯原子钟利用的两能级的能量差量级为 $10^{-5} \UeV$，跃迁发射的光子的频率量级为（普朗克常量 $h = 6.63\times10^{-34} \UJcdots$，元电荷 $e = 1.60\times10^{-19} \UC$）\xzanswer{C}

\fourchoices[answer=C]
{$10^{3} \UHz$}
{$10^{6} \UHz$}
{$10^{9} \UHz$}
{$10^{12} \UHz$}

%% answer: C
%% memo:
\memoanswer{
铯原子利用的两能极的能量差量级对应的能量为

$\varepsilon=10^{-5}\UeV=10^{-5}\times1.6\times10^{-19}\UJ=1.6\times10^{-24}\UJ$

由光子能量的表达式 $E = h\nu$ 可得，跃迁发射的光子的频率量级为

$\nu=\frac{\varepsilon}{h}=\frac{1.6\times10^{-24}}{6.63\times10^{-34}}\UHz\approx2.4\times10^{9}\UHz$

跃迁发射的光子的频率量级为 $10^{9} \UHz$。故选 C。
}

\item
%% number: 4
%% paperName: 2023年全国高考新课标第 4 题 6 分
%% typeId: 1
%% type: 单选题
%% chapter: 曲线运动
%% point: 人造地球卫星
%% method: 无
%% score: 6
%% degree: 750
%% duplicateId: 0
%% body:
2023 年 5 月，世界现役运输能力最大的货运飞船天舟六号，携带约 $5\,800\Ukg$ 的物资进入距离地面约 $400\Ukm$（小于地球同步卫星与地面的距离）的轨道，顺利对接中国空间站后近似做匀速圆周运动。对接后，这批物资\xzanswer{D}

\fourchoices[answer=D]
{质量比静止在地面上时小}
{所受合力比静止在地面上时小}
{所受地球引力比静止在地面上时大}
{做圆周运动的角速度大小比地球自转角速度大}

%% answer: D
%% memo:
\memoanswer{
A．物体在低速（速度远小于光速）宏观条件下质量保持不变，即在空间站和地面质量相同，故 A 错误；

BC．设空间站离地面的高度为 $h$，这批物质在地面上静止合力为零，在空间站所受合力为万有引力即

$F=\frac{GMm}{\left(R+h\right)^{2}}$

在地面受地球引力为

$F_{1}=\frac{GMm}{R^{2}}$

因此有 $F_{1} > F$，故 BC 错误；

D．物体绕地球做匀速圆周运动万有引力提供向心力

$\frac{GMm}{r^{2}}=m\omega^{2}r$

解得

$\omega=\sqrt{\frac{GM}{r^{3}}}$

这批物质在空间站内的轨道半径小于同步卫星的轨道半径，因此这批物质的角速度大于同步卫星的角速度，同步卫星的角速度等于地球自转的角速度，即这批物质的角速度大于地球自转的角速度，故 D 正确。

故选 D。
}

\item
%% number: 5
%% paperName: 2023年全国高考新课标第 5 题 6 分
%% typeId: 1
%% type: 单选题
%% chapter: 磁场
%% point: 带电粒子在复合场中的运动
%% method: 无
%% score: 6
%% degree: 650
%% duplicateId: 0
%% body:
一电子和一 $\alpha$ 粒子从铅盒上的小孔 O 竖直向上射出后，打到铅盒上方水平放置的屏幕 P 上的 a 和 b 两点，a 点在小孔 O 的正上方，b 点在 a 点的右侧，如图所示。已知 $\alpha$ 粒子的速度约为电子速度的 $\frac{1}{10}$，铅盒与屏幕之间存在匀强电场和匀强磁场，则电场和磁场方向可能为\xzanswer{C}

\onepicture[w=3.0cm]{figs/05.png}

\fourchoices[answer=C]
{电场方向水平向左、磁场方向垂直纸面向里}
{电场方向水平向左、磁场方向垂直纸面向外}
{电场方向水平向右、磁场方向垂直纸面向里}
{电场方向水平向右、磁场方向垂直纸面向外}

%% answer: C
%% memo:
\memoanswer{
A．带电粒子在电场和磁场中运动，打到 a 点的粒子电场力和洛伦兹力平衡，当电场向左磁场垂直直面向里时，$\alpha$ 粒子受到向左的电场力和洛伦兹力，电子受到向右的电场力和洛伦兹力均不能满足受力平衡打到 a 点，A 错误；

B．电场方向向左，磁场方向向外，此时如果 $\alpha$ 粒子打在 a 点则受到向左的电场力和向右的洛伦兹力平衡 $qE = qvB$，则电子向左的洛伦兹力大于向右的电场力向左偏转，同理如果电子打在 a 点，则 $\alpha$ 粒子向左的电场力大于向右的洛伦兹力向左偏转，均不会打在 b 点，B 错误；

CD．电场方向向右，磁场垂直纸面向里，如果 $\alpha$ 粒子打在 a 点，即向右的电场力和向左的洛伦兹力平衡 $qE = qvB$，电子向右的洛伦兹力大于向左的电场力向右偏转，同理如果电子打在 a，则 $\alpha$ 粒子向右的电场力大于向左的洛伦兹力向右偏转，均会打在 b 点；同理电场向右磁场垂直纸面向外时，$\alpha$ 粒子受到向右的电场力和洛伦兹力，电子受到向左的电场力和洛伦兹力不能受力平衡打到 a 点，故 C 正确 D 错误；

故选 C。
}

\item
%% number: 6
%% paperName: 2023年全国高考新课标第 6 题 6 分
%% typeId: 2
%% type: 多选题
%% chapter: 运动和力
%% point: 牛顿定律的应用
%% method: 无
%% score: 6
%% degree: 650
%% duplicateId: 0
%% body:
使甲、乙两条形磁铁隔开一段距离，静止于水平桌面上，甲的 N 极正对着乙的 S 极，甲的质量大于乙的质量，两者与桌面之间的动摩擦因数相等。现同时释放甲和乙，在它们相互接近过程中的任一时刻\xzanswer{BD}

\onepicture[w=9.5cm]{figs/06.png}

\fourchoices[answer=BD]
{甲的速度大小比乙的大}
{甲的动量大小比乙的小}
{甲的动量大小与乙的相等}
{甲和乙的动量之和不为零}

%% answer: BD
%% memo:
\memoanswer{
对甲、乙两条形磁铁分别做受力分析，如图所示

\onepicture[w=9.5cm]{figs/06_an.png}

A．根据牛顿第二定律有 $a_{\text{甲}} = \frac{F - \mu m_{\text{甲}}g}{m_{\text{甲}}}$，$a_{\text{乙}} = \frac{F - \mu m_{\text{乙}}g}{m_{\text{乙}}}$

由于 $m_{\text{甲}} > m_{\text{乙}}$，所以 $a_{\text{甲}} < a_{\text{乙}}$，由于两物体运动时间相同，且同时由静止释放，可得 $v_{\text{甲}} < v_{\text{乙}}$

A 错误；

BCD．对于整个系统而言，由于 $\mu m_{\text{甲}}g > \mu m_{\text{乙}}g$，合力方向向左，合冲量方向向左，所以合动量方向向左，显然甲的动量大小比乙的小，BD 正确、C 错误。

故选 BD。
}

\item
%% number: 7
%% paperName: 2023年全国高考新课标第 7 题 6 分
%% typeId: 2
%% type: 多选题
%% chapter: 功和能
%% point: 功能中的图象
%% method: 无
%% score: 6
%% degree: 550
%% duplicateId: 0
%% body:
一质量为 $1 \Ukg$ 的物体在水平拉力的作用下，由静止开始在水平地面上沿 $x$ 轴运动，出发点为 $x$ 轴零点，拉力做的功 $W$ 与物体坐标 $x$ 的关系如图所示。物体与水平地面间的动摩擦因数为 0.4，重力加速度大小取 $10 \Umsq$。下列说法正确的是\xzanswer{BC}

\onepicture[w=6.2cm]{figs/07.png}

\fourchoices[answer=BC]
{在 $x = 1 \Um$ 时，拉力的功率为 $6 \UW$}
{在 $x = 4 \Um$ 时，物体的动能为 $2 \UJ$}
{从 $x = 0$ 运动到 $x = 2 \Um$，物体克服摩擦力做的功为 $8 \UJ$}
{从 $x = 0$ 运动到 $x = 4 \Um$ 的过程中，物体的动量最大为 $2 \Ukgms$}

%% answer: BC
%% memo:
\memoanswer{
由于拉力在水平方向，则拉力做的功为 $W = Fx$

可看出 $W - x$ 图像的斜率代表拉力 $F$。

AB．在物体运动的过程中根据动能定理有

$W - \mu mgx = \frac{1}{2}mv^{2}$

则 $x = 1 \Um$ 时物体的速度为 $v_{1} = 2 \Ums$

$x = 1 \Um$ 时，拉力为 $F = \frac{\Delta W}{\Delta x} = 6 \UN$

则此时拉力的功率 $P = Fv_{1} = 12 \UW$

$x = 4 \Um$ 时物体的动能为 $E_{k} = 2 \UJ$

A 错误、B 正确；

C．从 $x = 0$ 运动到 $x = 2 \Um$，物体克服摩擦力做的功为 $W_{f} = \mu mgx = 8 \UJ$。C 正确；

D．根据 $W - x$ 图像可知在 $0 \sim 2 \Um$ 的过程中 $F_{1} = 6 \UN$，$2 \sim 4 \Um$ 的过程中 $F_{2} = 3 \UN$，由于物体受到的摩擦力恒为 $f = 4 \UN$，则物体在 $x = 2 \Um$ 处速度最大，且根据选项 AB 分析可知此时的速度 $v_{2} = \sqrt{8} \Ums$。则从 $x = 0$ 运动到 $x = 4 \Um$ 的过程中，物体的动量最大为 $p = mv = 2\sqrt{2} \Ukgms$。D 错误。

故选 BC。
}

\item
%% number: 8
%% paperName: 2023年全国高考新课标第 8 题 6 分
%% typeId: 2
%% type: 多选题
%% chapter: 气体性质
%% point: 气体连接体
%% method: 无
%% score: 6
%% degree: 450
%% duplicateId: 0
%% body:
如图，一封闭着理想气体的绝热汽缸置于水平地面上，用轻弹簧连接的两绝热活塞将汽缸分为 f、g、h 三部分，活塞与汽缸壁间没有摩擦。初始时弹簧处于原长，三部分中气体的温度、体积、压强均相等。现通过电阻丝对 f 中的气体缓慢加热，停止加热并达到稳定后\xzanswer{AD}

\onepicture[w=7.5cm]{figs/08.png}

\fourchoices[answer=AD]
{h 中的气体内能增加}
{f 与 g 中的气体温度相等}
{f 与 h 中的气体温度相等}
{f 与 h 中的气体压强相等}

%% answer: AD
%% memo:
\memoanswer{
A．当电阻丝对 f 中的气体缓慢加热时，f 中的气体内能增大，温度升高，根据理想气体状态方程可知 f 中的气体压强增大，会缓慢推动左边活塞，则弹簧被压缩。与此同时弹簧对右边活塞有弹力作用，缓慢向右推动左边活塞。故活塞对 h 中的气体做正功，且是绝热过程，由热力学第一定律可知，h 中的气体内能增加，A 正确；

B．未加热前，三部分中气体的温度、体积、压强均相等，当系统稳定时，活塞受力平衡，可知弹簧处于压缩状态，对左边活塞分析

$p_{f}S=F_{\text{弹}}+p_{g}S$

则

$p_{f}>p_{g}$

分别对 f、g 内的气体分析，根据理想气体状态方程有

$\frac{p_{0}V_{0}}{T_{0}}=\frac{p_{f}V_{f}}{T_{f}}$

$\frac{p_{0}V_{0}}{T_{0}}=\frac{p_{g}V_{g}}{T_{g}}$

由题意可知，因弹簧被压缩，则 $V_{f} > V_{g}$，联立可得 $T_{f} > T_{g}$。B 错误；

C．在达到稳定过程中 h 中的气体体积变小，压强变大，f 中的气体体积变大。由于稳定时弹簧保持平衡状态，故稳定时 f、h 中的气体压强相等，根据理想气体状态方程对 h 气体分析可知

$\frac{p_{0}V_{0}}{T_{0}}=\frac{p_{h}V_{h}}{T_{h}}$

联立可得 $T_{f} > T_{h}$。C 错误；

D．对弹簧、活塞及 g 中的气体组成的系统分析，根据平衡条件可知，f 与 h 中的气体压强相等，D 正确。

故选 AD。
}

\item
%% number: 9
%% paperName: 2023年全国高考新课标第 9 题 9 分
%% typeId: 4
%% type: 实验题
%% chapter: 电路
%% point: 电容器充放电实验
%% method: 无
%% score: 9
%% degree: 550
%% duplicateId: 0
%% body:
在“观察电容器的充、放电现象”实验中，所用器材如下：电池、电容器、电阻箱、定值电阻、小灯泡、多用电表、电流表、秒表、单刀双掷开关以及导线若干。
\begin{enumerate}
	\item
	用多用电表的电压挡检测电池的电压。检测时，红表笔应该与电池的\tkanswer{正极}（填“正极”或“负极”）接触。
	\item
	某同学设计的实验电路如图 \ref{2023全国新课标09a} 所示。先将电阻箱的阻值调为 $R_{1}$，将单刀双掷开关 S 与“1”端相接，记录电流随时间的变化。电容器充电完成后，开关 S 再与“2”端相接，相接后小灯泡亮度变化情况可能是\tkanswer{C}。（填正确答案标号）

	A．迅速变亮，然后亮度趋于稳定

	B．亮度逐渐增大，然后趋于稳定

	C．迅速变亮，然后亮度逐渐减小至熄灭
	\item
	将电阻箱的阻值调为 $R_{2}$（$R_{2} > R_{1}$），再次将开关 S 与“1”端相接，再次记录电流随时间的变化情况。两次得到的电流 $I$ 随时间 $t$ 变化如图 \ref{2023全国新课标09b} 中曲线所示，其中实线是电阻箱阻值为\tkanswer{$R_{2}$}（填“$R_{1}$”或“$R_{2}$”）时的结果，曲线与坐标轴所围面积等于该次充电完成后电容器上的\tkanswer{电荷量}（填“电压”或“电荷量”）。
\end{enumerate}
\twopicture[labelA=2023全国新课标09a, labelB=2023全国新课标09b, wA=4.2cm, wB=5.2cm, align=right]
{figs/09a.png}{figs/09b.png}

%% answer:
\jdanswer{
\begin{enumerate}
	\item
	正极

	\item
	C

	\item
	$R_{2}$，电荷量

\end{enumerate}
}
%% memo:
\memoanswer{
（1）多用电表红表笔流入电流，黑表笔流出电流，故电流表红表笔应该与电池的正极接触；

（2）电容器充电完成后，开始时两极板电量较多，电势差较大，当闭合“2”接入小灯泡，回路立即形成电流，灯泡的迅速变亮；随着时间的积累，两极板电量变少，电势差变小，流过灯泡的电流减小，直至两极板电荷量为零不带电，则无电流流过小灯泡即熄灭，故选 C。

（3）开始充电时两极板的不带电，两极板电势差为零，设电源内阻为 $r$，则开始充电时有 $E = I(R + r)$，由图像可知开始充电时实线的电流较小，故电路中的电阻较大，因此电阻箱阻值为 $R_{2}$；

图像的物理意义为充电过程中电流随时间的变化图线，故曲线与坐标轴所围面积等于该次充电完成后电容器上的电荷量。
}

\item
%% number: 10
%% paperName: 2023年全国高考新课标第 10 题 12 分
%% typeId: 4
%% type: 实验题
%% chapter: 机械振动
%% point: 单摆测重力加速度
%% method: 无
%% score: 12
%% degree: 450
%% duplicateId: 0
%% body:
一学生小组做“用单摆测量重力加速度的大小”实验。
\begin{enumerate}
	\item
	用实验室提供的螺旋测微器测量摆球直径。首先，调节螺旋测微器，拧动微调旋钮使测微螺杆和测砧相触时，发现固定刻度的横线与可动刻度上的零刻度线未对齐，如图 \ref{2023全国新课标10a} 所示，该示数为\tkanswer{0.007}$\Umm$；螺旋测微器在夹有摆球时示数如图 \ref{2023全国新课标10b} 所示，该示数为\tkanswer{20.034}$\Umm$，则摆球的直径为\tkanswer{20.027}$\Umm$。
	\item
	单摆实验的装置示意图如图 \ref{2023全国新课标10c} 所示，其中角度盘需要固定在杆上的确定点 O 处，摆线在角度盘上所指的示数为摆角的大小。若将角度盘固定在 O 点上方，则摆线在角度盘上所指的示数为 $5\Udu$ 时，实际摆角\tkanswer{大于}$5\Udu$（填“大于”或“小于”）。
	\item
	某次实验所用单摆的摆线长度为 $81.50\Ucm$，则摆长为\tkanswer{82.5}$\Ucm$。实验中观测到从摆球第 1 次经过最低点到第 61 次经过最低点的时间间隔为 $54.60\Us$，则此单摆周期为\tkanswer{1.82}$\Us$，该小组测得的重力加速度大小为\tkanswer{9.83}$\Umsq$（结果均保留 3 位有效数字，$\pi^{2}$ 取 9.870）
\end{enumerate}
\threepicture[labelA=2023全国新课标10a, labelB=2023全国新课标10b, labelC=2023全国新课标10c, wA=3.0cm, wB=5.3cm, wC=2.4cm, align=right]
{figs/10a.png}{figs/10b.png}{figs/10c.png}

%% answer:
\jdanswer{
\begin{enumerate}
	\item
	0.007，20.034，20.027

	\item
	大于

	\item
	82.5，1.82，9.83

\end{enumerate}
}
%% memo:
\memoanswer{
（1）测量前测微螺杆与和测砧相触时，图（a）的示数为

$d_{0}=0\Umm+0.7\times0.01\Umm=0.007\Umm$

螺旋测微器读数是固定刻度读数（0.5$\Umm$ 的整数倍）加可动刻度（0.5$\Umm$ 以下的小数）读数，图中读数为

$d_{1}=20\Umm+3.4\times0.01\Umm=20.034\Umm$

则摆球的直径为

$d=d_{1}-d_{0}=20.027\Umm$

（2）角度盘的大小一定，即在规定的位置安装角度盘，测量的摆角准确，但将角度盘固定在规定位置上方，即角度盘到悬挂点的距离变短，同样的角度，摆线在刻度盘上扫过的弧长变短，故摆线在角度盘上所指的示数为 $5\Udu$ 时，实际摆角大于 $5\Udu$；

（3）单摆的摆线长度为 $81.50\Ucm$，则摆长为

$l=l_{0}+\frac{d}{2}=81.50\Ucm+\frac{2.0027}{2}\Ucm=82.5\Ucm$

结果保留三位有效数字，得摆长为 $82.5\Ucm$；

一次全振动单摆经过最低点两次，故此单摆的周期为

$T=\frac{2t}{N}=\frac{54.60}{30}\Us=1.82\Us$

由单摆的周期表达式 $T = 2\pi \sqrt{\frac{l}{g}}$ 得，重力加速度

$g = \frac{4\pi^{2}l}{T^{2}} = 9.83\Umsq$
}

\item
%% number: 11
%% paperName: 2023年全国高考新课标第 11 题 10 分
%% typeId: 6
%% type: 计算题
%% chapter: 曲线运动
%% point: 平抛运动
%% method: 无
%% score: 10
%% degree: 650
%% duplicateId: 0
%% body:
将扁平的石子向水面快速抛出，石子可能会在水面上一跳一跳地飞向远方，俗称“打水漂”。要使石子从水面跳起产生“水漂”效果，石子接触水面时的速度方向与水面的夹角不能大于 $\theta$。为了观察到“水漂”，一同学将一石子从距水面高度为 $h$ 处水平抛出，抛出速度的最小值为多少？（不计石子在空中飞行时的空气阻力，重力加速度大小为 $g$）

%% answer:
\jdanswer{
$\frac{\sqrt{2gh}}{\tan \theta}$
}
%% memo:
\memoanswer{
石子做平抛运动，竖直方向做自由落体运动，则有

$2gh = v_{y}^{2}$

可得落到水面上时的竖直速度 $v_{y} = \sqrt{2gh}$

由题意可知 $\frac{v_{y}}{v_{0}} \leq \tan \theta$

即 $v_{0} \geq \frac{\sqrt{2gh}}{\tan \theta}$

石子抛出速度的最小值为 $\frac{\sqrt{2gh}}{\tan \theta}$。
}

\item
%% number: 12
%% paperName: 2023年全国高考新课标第 12 题 14 分
%% typeId: 6
%% type: 计算题
%% chapter: 电场
%% point: 带电粒子的直线运动
%% method: 无
%% score: 14
%% degree: 550
%% duplicateId: 0
%% body:
密立根油滴实验的示意图如图所示。两水平金属平板上下放置，间距固定，可从上板中央的小孔向两板间喷入大小不同、带电量不同、密度相同的小油滴。两板间不加电压时，油滴 a、b 在重力和空气阻力的作用下竖直向下匀速运动，速率分别为 $v_{0}$、$\frac{v_{0}}{4}$；两板间加上电压后（上板为正极），这两个油滴很快达到相同的速率 $\frac{v_{0}}{2}$，均竖直向下匀速运动。油滴可视为球形，所受空气阻力大小与油滴半径、运动速率成正比，比例系数视为常数。不计空气浮力和油滴间的相互作用。
\begin{enumerate}
	\item
	求油滴 a 和油滴 b 的质量之比；
	\item
	判断油滴 a 和油滴 b 所带电荷的正负，并求 a、b 所带电荷量的绝对值之比。
\end{enumerate}
\onepicture[w=6.5cm, align=right]{figs/12.png}

%% answer:
\jdanswer{
\begin{enumerate}
	\item
	$8:1$

	\item
	油滴 a 带负电，油滴 b 带正电；$4:1$

\end{enumerate}
}
%% memo:
\memoanswer{
（1）设油滴半径 $r$，密度为 $\rho$，则油滴质量

$m=\frac{4}{3}\pi r^{3}\rho$

则速率为 $v$ 时受阻力

$f = krv$

则当油滴匀速下落时

$mg = f$

解得

$r=\sqrt{\frac{3kv}{4\pi\rho g}}\propto\sqrt{v}$

可知

$\frac{r_{a}}{r_{b}}=\sqrt{\frac{v_{0}}{\frac{1}{4}v_{0}}}=2$

则

$\frac{m_{a}}{m_{b}}=\frac{r_{a}^{3}}{r_{b}^{3}}=\frac{8}{1}$

（2）两板间加上电压后（上板为正极），这两个油滴很快达到相同的速率 $\frac{v_{0}}{2}$，可知油滴 a 做减速运动，油滴 b 做加速运动，可知油滴 a 带负电，油滴 b 带正电；当再次匀速下落时，对 a 由受力平衡可得

$\left|q_{a}\right|E+f_{a}=m_{a}g$

其中

$f_{a}=\frac{\frac{v_{0}}{2}}{v_{0}}m_{a}g=\frac{1}{2}m_{a}g$

对 b 由受力平衡可得

$f_{b}-q_{b}E=m_{b}g$

其中

$f_{b}=\frac{\frac{v_{0}}{2}}{\frac{1}{4}v_{0}}m_{b}g=2m_{b}g$

联立解得

$\left|\frac{q_{a}}{q_{b}}\right|=\frac{m_{a}}{2m_{b}}=\frac{4}{1}$
}

\item
%% number: 13
%% paperName: 2023年全国高考新课标第 13 题 20 分
%% typeId: 6
%% type: 计算题
%% chapter: 电磁感应
%% point: 电磁感应能量分析
%% method: 无
%% score: 20
%% degree: 350
%% duplicateId: 0
%% body:
一边长为 $L$、质量为 $m$ 的正方形金属细框，每边电阻为 $R_{0}$，置于光滑的绝缘水平桌面（纸面）上。宽度为 $2L$ 的区域内存在方向垂直于纸面的匀强磁场，磁感应强度大小为 $B$，两虚线为磁场边界，如图 \ref{2023全国新课标13a} 所示。
\begin{enumerate}
	\item
	使金属框以一定的初速度向右运动，进入磁场。运动过程中金属框的左、右边框始终与磁场边界平行，金属框完全穿过磁场区域后，速度大小降为它初速度的一半，求金属框的初速度大小。
	\item
	在桌面上固定两条光滑长直金属导轨，导轨与磁场边界垂直，左端连接电阻 $R_{1} = 2R_{0}$，导轨电阻可忽略，金属框置于导轨上，如图 \ref{2023全国新课标13b} 所示。让金属框以与（1）中相同的初速度向右运动，进入磁场。运动过程中金属框的上、下边框处处与导轨始终接触良好。求在金属框整个运动过程中，电阻 $R_{1}$ 产生的热量。
\end{enumerate}
\twopicture[labelA=2023全国新课标13a, labelB=2023全国新课标13b, wA=4.6cm, wB=7.6cm, align=right]
{figs/13a.png}{figs/13b.png}

%% answer:
\jdanswer{
\begin{enumerate}
	\item
	$\frac{B^{2}L^{3}}{mR_{0}}$

	\item
	$\frac{3B^{4}L^{6}}{25mR_{0}^{2}}$

\end{enumerate}
}
%% memo:
\memoanswer{
（1）金属框进入磁场过程中有

$\overline{E}=BL\frac{L}{t}$

则金属框进入磁场过程中流过回路的电荷量为

$q_{1}=\frac{\overline{E}}{4R_{0}}t=\frac{BL^{2}}{4R_{0}}$

则金属框完全穿过磁场区域的过程中流过回路的电荷量为

$q=\frac{BL^{2}}{2R_{0}}$

且有

$-BqL=\frac{mv_{0}}{2}-mv_{0}$

联立有 $v_{0}=\frac{B^{2}L^{3}}{mR_{0}}$

（2）设金属框的初速度为 $v_{0}$，则金属框进入磁场时的末速度为 $v_{1}$，向右为正方向。由于导轨电阻可忽略，此时金属框上下部分被短路，故电路中的总电阻

$R_{\text{总}}=R_{0}+\frac{2R_{0}\cdot R_{0}}{2R_{0}+R_{0}}=\frac{5R_{0}}{3}$

再根据动量定理有

$-\frac{B^{2}L^{3}}{R_{\text{总}}}=mv_{1}-mv_{0}$

解得

$v_{1}=\frac{2B^{2}L^{3}}{5mR_{0}}$

则在此过程中根据能量守恒有

$\frac{1}{2}mv_{0}^{2}=Q_{1}+\frac{1}{2}mv_{1}^{2}$

解得

$Q_{1}=\frac{21B^{4}L^{6}}{50mR_{0}^{2}}$

其中

$Q_{R_{1}}=\frac{2}{15}Q_{1}=\frac{7B^{4}L^{6}}{125mR_{0}^{2}}$

此后线框完全进入磁场中，则线框左右两边均作为电源，且等效电路图如下

\onepicture[w=5.5cm]{figs/13_an.png}

则此时回路的总电阻

$R_{\text{总}}'=2R_{0}+\frac{R_{0}}{2}=\frac{5R_{0}}{2}$

设线框刚离开磁场时的速度为 $v_{2}$，再根据动量定理有

$-\frac{B^{2}L^{3}}{R_{\text{总}}'}=mv_{2}-mv_{1}$

解得 $v_{2}=0$

则说明线框刚离开磁场时就停止运动了，则再根据能量守恒有

$\frac{1}{2}mv_{1}^{2}=Q_{2}$

其中

$Q_{R_{1}}'=\frac{4}{5}Q_{2}=\frac{8B^{4}L^{6}}{125mR_{0}^{2}}$

则在金属框整个运动过程中，电阻 $R_{1}$ 产生的热量

$Q_{R1\text{总}}=Q_{R1}+Q_{R1}'=\frac{3B^{4}L^{6}}{25mR_{0}^{2}}$
}

\end{enumerate}

\end{document}
